\section{Sobolev capacity and two Hessian factors}
\label{sec:hessian-sobolev-two-factors}

Weak convergence of measures is not enough to pass to the limit in
their products with bounded potentials; one needs a capacity that
dominates the measures and makes the potentials quasicontinuous.
Sobolev capacity does this for the first Hessian of a bounded
admissible potential, through its linear trace, but a product of two
Hessians can charge a set of zero Sobolev capacity. The trace is taken
against a $\ddc$-closed conformal multiple of $\omega^{n-1}$: an
admissible form need not be positive in every tangent direction, so
only its trace relative to $\omega$ is known to be nonnegative.

Choose $\rho>0$ by \cref{thm:gauduchon} for $\omega$ and put
\[
 T=\rho\omega^{n-1},\qquad \ddc T=0,\qquad
 A=\int_X\omega\wedge T>0.
\]
Write $H^1(X)=W^{1,2}(X,\mathbb R)$ and use the equivalent Sobolev norm
\[
 \|f\|_{H^1}^2=\int_Xf^2\omega^n+
                   \int_X\mathrm df\wedge\dc f\wedge T.
\]
For an open set $G\subset X$ define
\[
 \operatorname{Cap}_1(G)
 =\inf\{\|f\|_{H^1}^2:f\in H^1(X),\ f\geq1
                          \text{ almost everywhere on }G\},
\]
and for an arbitrary set $E$ take the infimum over open $G\supset E$.
It is distinct from the local Hessian capacities of
\cref{sec:hessian-smooth-capacity}. Truncation allows $0\leq f\leq1$ in the infimum for open sets.
\index{Capacity!Sobolev}

\begin{lemma}\label{hessian:sobolev-capacity-tools}
The capacity $\operatorname{Cap}_1$ is countably subadditive. For a
smooth real $f$ and $a>0$,
\[
 \operatorname{Cap}_1(\{|f|>a\})\leq a^{-2}\|f\|_{H^1}^2.
\]
If a positive Radon measure $\mu$ satisfies
$|\int_X\phi\,\mathrm d\mu|\leq C\|\phi\|_{H^1}$ for every smooth real
$\phi$, then for every Borel set $E$,
\[
 \mu(E)\leq C\operatorname{Cap}_1(E)^{1/2}.
\]
\end{lemma}
\begin{proof}
Let $G_i$ be open sets with $\sum_i\operatorname{Cap}_1(G_i)<\infty$;
if the sum is infinite there is nothing to prove. Fix $\epsilon>0$ and
choose Sobolev functions $0\leq f_i\leq1$ equal to one almost
everywhere on $G_i$, with
$\|f_i\|_{H^1}^2\leq\operatorname{Cap}_1(G_i)+\epsilon2^{-i}$.
The Sobolev chain rule gives
\[
 \|\max_{i\leq N}f_i\|_{H^1}^2\leq\sum_{i\leq N}\|f_i\|_{H^1}^2.
\]
Indeed, almost everywhere the gradient of the maximum is the gradient
of an active function, and the squared value is bounded by the sum of
the squared values. These maxima increase with $N$, converge in
$L^2$, are bounded in $H^1$, and have a weak $H^1$ limit equal to one
almost everywhere on $\bigcup_iG_i$. Lower semicontinuity of the norm
gives
$\operatorname{Cap}_1(\bigcup_iG_i)\leq\sum_i\operatorname{Cap}_1(G_i)+\epsilon$.
This proves subadditivity for open sets, and taking open supersets
proves it for arbitrary sets.
The level-set estimate follows by testing with
$\min\{1,|f|/a\}$ and the same Sobolev chain rule.

Fix a compact $K\subset G$ and a Sobolev competitor $0\leq f\leq1$
for the open set $G$. Mollify $f$ in finitely many coordinate charts
by nonnegative kernels and glue the results with a fixed partition of
unity. This gives smooth nonnegative $f_\epsilon\to f$ in $H^1$ that
satisfy $f_\epsilon=1$ on $K$ for small $\epsilon$: every convolution
ball centered on $K$ then lies in $G$, where $f=1$ almost everywhere.
The assumed estimate gives
$\mu(K)\leq C\lim_\epsilon\|f_\epsilon\|_{H^1}=C\|f\|_{H^1}$.
Take the infimum over $f$ and then the supremum over compact $K\subset G$
to obtain the estimate for $G$. For a Borel set $E$, apply it to
every open $G\supset E$ and use $\mu(E)\leq\mu(G)$.
\end{proof}

\begin{lemma}\label{hessian:bounded-sobolev-trace}\index{Sobolev regularity!bounded Hessian potentials}
Every bounded admissible $u$ belongs to $H^1(X)$. Its trace
\[
 \mu_u=(\omega+\ddc u)\wedge T
\]
is a positive Radon measure of total mass $A$. If $\|u\|_\infty\leq M$,
then
\[
 \|u\|_{H^1}\leq C(1+M),\qquad
 \mu_u(E)\leq C(1+M)\operatorname{Cap}_1(E)^{1/2}
\]
for every Borel $E$. On each fixed relatively compact coordinate chart,
the total variation of every coefficient of $\ddc u$ is also bounded
on $E$ by $C(1+M)\operatorname{Cap}_1(E)^{1/2}$.
\end{lemma}
\begin{proof}
First let $u$ be smooth admissible. The first-cone inequality gives
$\mu_u\geq0$, since $T$ is conformal to $\omega^{n-1}$, and
$\ddc T=0$ gives $\mu_u(X)=A$. Stokes' theorem and $\ddc T=0$ give
$\int_X\ddc(u^2)\wedge T=0$. Since
$\ddc(u^2)=2u\ddc u+2\mathrm du\wedge\dc u$, this yields the energy
identity
\[
 \int_X\mathrm du\wedge\dc u\wedge T
   =-\int_Xu\,\ddc u\wedge T
   =-\int_Xu\,\mathrm d\mu_u+\int_Xu\omega\wedge T
   \leq2MA.
\]
Together with the bound $M^2\int_X\omega^n$ for the $L^2$ term, this
gives $\|u\|_{H^1}\leq C(1+M)$.
For bounded admissible $u$, choose smooth admissible $u_j\downarrow u$
by \cref{hessian:bounded-decreasing-approximation}, discard a finite
initial segment, and take $\|u_j\|_\infty\leq M+1$ as in
\cref{hessian:bounded-zero-order}. Bounded convergence gives strong
$L^2$ convergence. The energy bound gives weak $H^1$ compactness, and
every weak cluster point agrees with $u$ in $L^2$; thus the entire
sequence converges weakly to $u$ in $H^1$. Lower semicontinuity of the
norm proves the Sobolev bound.
The trace measures converge weakly to $\mu_u$, preserving positivity and total mass.

For any smooth real $\phi$, the first-order Stokes identity is
\[
 \int_X\phi\,\ddc u\wedge T
  =-\int_X\mathrm d\phi\wedge\dc u\wedge T
       +\int_X\phi\,\dc u\wedge\mathrm dT.
\]
It remains valid for $u\in H^1$ by smooth Sobolev approximation.
Both terms are bounded by $C\|u\|_{H^1}\|\phi\|_{H^1}$; the background term
$\int\phi\omega\wedge T$ is bounded by $C\|\phi\|_{H^1}$.
Hence $|\int_X\phi\,\mathrm d\mu_u|\leq C(1+M)\|\phi\|_{H^1}$, and
\cref{hessian:sobolev-capacity-tools} gives the bound for $\mu_u$.

For a smooth admissible approximant, the eigenvalues of
$\omega+\ddc u_j$ relative to $\omega$ have modulus at most its trace,
by the proof of \cref{hessian:weighted-cone-trace}, since $m\geq2$.
Consequently every coefficient of this form is dominated by a fixed
multiple of its trace.
On a fixed coordinate chart this bounds the coefficient variations of
$\ddc u_j$ by $C(\mu_{u_j}+\omega^n)$. Passing to the weak measure
limits of \cref{hessian:bounded-zero-order} gives the same domination
by $C(\mu_u+\omega^n)$. Finally
$\int_E\omega^n\leq\operatorname{Cap}_1(E)$ and
$\operatorname{Cap}_1(E)\leq\|1\|_{H^1}^2$ absorb the smooth term into the
claimed square-root estimate.
\end{proof}

\begin{lemma}\label{hessian:bounded-sobolev-quasicontinuity}\index{Admissible potential!Sobolev quasicontinuity}
Every bounded admissible potential is continuous on the
complement of an open set of arbitrarily small $\operatorname{Cap}_1$.
The same open set can be chosen for any countable family of such
potentials.
\end{lemma}
\begin{proof}
Choose uniformly bounded smooth admissible $u_j\downarrow u$ by
\cref{hessian:bounded-decreasing-approximation}. By the preceding lemma
these functions converge weakly to $u$ in $H^1$. Mazur's lemma gives
finite convex combinations $w_N$ of terms with $j\geq N$ such that
\[
 \|w_N-u\|_{H^1}\leq8^{-N}.
\]
Because $u\leq w_N\leq u_N$, they converge pointwise to the canonical
function $u$, not just to an unspecified almost-everywhere representative.
Put
\[
 G_{N_0}=\bigcup_{N\geq N_0}
                  \{|w_{N+1}-w_N|>2^{-N}\}.
\]
\Cref{hessian:sobolev-capacity-tools} implies
\[
 \operatorname{Cap}_1(G_{N_0})
 \leq\sum_{N\geq N_0}4^N\|w_{N+1}-w_N\|_{H^1}^2
 \leq C\sum_{N\geq N_0}16^{-N}\longrightarrow0.
\]
Outside $G_{N_0}$, the successive differences form a uniformly summable
series. Thus $w_N$ converge uniformly there and their limit $u$ is
continuous on that closed set. For a countable family, choose the
exceptional open sets with capacities less than $\epsilon2^{-i}$ and
take their union. Countable subadditivity proves the last assertion.
\end{proof}

The capacity estimate depends on using a first Hessian. A product
of two Hessians can already charge a set of zero Sobolev capacity.

\begin{lemma}\label{hessian:sobolev-codimension-two}\index{Capacity!Sobolev zero sets}
Let $E$ be a compact subset of a smooth real submanifold of $X$ of
codimension two. Then $\operatorname{Cap}_1(E)=0$.
\end{lemma}
\begin{proof}
It suffices to treat a compact subset of a submanifold chart, since
finitely many such subsets cover $E$ and the capacity is subadditive.
Choose real coordinates $(s,t)$ there, with $t\in\mathbb R^2$ and the
submanifold given by $t=0$. Choose $0\leq\eta\leq1$ smooth,
supported in this chart and equal to one near the compact subset.
Put $r=|t|$.
For $0<\epsilon<1/2$ define
\[
 h_\epsilon(r)=
 \begin{cases}
  1,&0\leq r\leq\epsilon^2,\\
  \displaystyle\frac{\log(\epsilon/r)}{\log(1/\epsilon)},
       &\epsilon^2<r<\epsilon,\\
  0,&r\geq\epsilon.
 \end{cases}
\]
The Sobolev function $f_\epsilon=\eta h_\epsilon$ equals one on an
open neighborhood of that subset. On the fixed compact coordinate
set the Sobolev norm is equivalent to the Euclidean one. Integration
in the two normal variables gives
\[
 \|f_\epsilon\|_{H^1}^2
 \leq C\left(\epsilon^2+
     \int_{\epsilon^2}^{\epsilon}
       \frac{\mathrm dr}{r\log(1/\epsilon)^2}\right)
 =C\left(\epsilon^2+\frac1{\log(1/\epsilon)}\right)
 \longrightarrow0.
\]
The derivatives of $\eta$ contribute at most $C\epsilon^2$.
Since $f_\epsilon$ is a competitor for an open neighborhood of the
subset, the capacity of the subset is at most
$\|f_\epsilon\|_{H^1}^2$ for every $\epsilon$, which proves the
assertion.
\end{proof}

\begin{example}\label{hessian:sobolev-two-product-obstruction}\index{Hessian measure!Sobolev capacity obstruction}
There are bounded continuous $\omega$-plurisubharmonic potentials
$u,v$ and a compact set $E$ such that
\[
 \operatorname{Cap}_1(E)=0,\qquad
 \int_E Q_2(u,v)\wedge\omega^{m-2}>0.
\]
Consequently the first-Hessian capacity estimate of
\cref{hessian:bounded-sobolev-trace} does not extend to
this positive two-factor measure, even within the continuous
$\omega$-plurisubharmonic subclass.
\end{example}
\begin{proof}
Fix a holomorphic chart with coordinates $z_j=x_j+\mathrm i y_j$ and a smooth
nonnegative cutoff $\chi$ supported in the chart and equal to one
on a smaller coordinate box. Extend
\[
 u=a\chi|x_1|,\qquad v=a\chi|x_2|
\]
by zero outside the chart. These functions are globally continuous
and bounded. To check their curvature, replace $|x_j|$ by
$g_\delta(x_j)=(x_j^2+\delta^2)^{1/2}$, with $0<\delta\leq1$.
The Hessian of $\chi g_\delta$ has the nonnegative term
$\chi\ddc g_\delta$; all its remaining terms are uniformly bounded
forms, since $g_\delta$ and $g_\delta'$ are uniformly bounded on the
fixed support. Thus $\ddc(\chi g_\delta)\geq-C\omega$ with $C$
independent of $\delta$. Choosing $a>0$ sufficiently small gives
$\omega+\ddc(a\chi g_\delta)\geq\omega/2$ for both coordinates.
Their decreasing limits $u,v$ as $\delta\downarrow0$ are therefore
$\omega$-plurisubharmonic and admissible.

On the box where $\chi=1$, set
$\alpha_j=\frac{\mathrm i}{4\pi}\mathrm dz_j\wedge\mathrm d\bar z_j$.
The book's normalization gives
\[
 \ddc g_\delta(x_j)=\tfrac12g_\delta''(x_j)\alpha_j,
 \qquad g_\delta''(x_j)
       =\frac{\delta^2}{(x_j^2+\delta^2)^{3/2}}
       \longrightarrow2\delta_{\{x_j=0\}}.
\]
The measures $g_\delta^{\prime\prime}(t)\,\mathrm dt$ have mass two,
while their mass outside any fixed neighborhood of zero tends to
zero. Thus $\Theta=\ddc u=a\delta_{\{x_1=0\}}\alpha_1$ on the box. This
current is $\mathrm d$- and $\dc$-closed and depends only on $x_1$,
while $v=a|x_2|$ there. Hence
$\ddc(g_\delta(x_2)\Theta)=\ddc g_\delta(x_2)\wedge\Theta$, and since
$g_\delta\to|\cdot|$ uniformly, \eqref{eq:hessian-bounded-recursion}
gives on the box
\[
 D_2(u,v)=\ddc(v\Theta)
 =\lim_{\delta\downarrow0}a\,\ddc g_\delta(x_2)\wedge\Theta
 =a^2\delta_{\{x_1=0\}}\delta_{\{x_2=0\}}
                         \alpha_1\wedge\alpha_2;
\]
the two factors depend on distinct real variables, so their product
converges to the tensor-product measure.
Choose $E$ to be a compact piece of $\{x_1=x_2=0\}$ inside this
box with nonempty relative interior. The preceding lemma gives
$\operatorname{Cap}_1(E)=0$. On this box, the positive measure
$D_2(u,v)\wedge\omega^{n-2}$ has strictly positive mass on $E$,
since $\omega$ is strictly positive. The smooth term $\omega^n$ and
the two single-Hessian terms in the expansion
\[
 Q_2(u,v)\wedge\omega^{m-2}
 =\omega^n+\omega^{n-1}\wedge(\ddc u+\ddc v)
                         +D_2(u,v)\wedge\omega^{n-2}
\]
have zero mass on $E$. Indeed, on the box $\omega^{n-1}\wedge\ddc u$
is a smooth multiple of the hypersurface measure of $\{x_1=0\}$,
which vanishes on $E$, and similarly for $v$. Hence the asserted
mass is strictly positive. A bound by
$C\operatorname{Cap}_1(F)^b$ for all Borel $F$, with any $b>0$,
would contradict this choice of $E$.
\end{proof}

The preceding Sobolev-capacity estimate therefore cannot be iterated
to control higher-degree coefficient measures. A local Hessian
capacity gives that control in the next section.
